% chapters/chapter3/3_6_tradeoff.tex
% !TeX root=../../main.tex

\section{جمع‌بندی نتایج}
\label{sec:results_tradeoff}

شبیه‌سازی‌ها و ارزیابی‌های انجام‌شده در این فصل، دیدگاهی جامع نسبت به چالش‌های هدایت سیستم‌های زیرفعال و نحوهٔ غلبه بر آن‌ها ارائه می‌دهد. مقایسهٔ رفتار کنترل‌کننده‌های خطی و غیرخطی در مواجهه با شرایط مختلف، منجر به شناسایی مجموعه‌ای از مصالحه‌های مهندسی (\lr{Trade-off}) در طراحی سیستم‌های کنترل جرثقیل سقفی می‌شود. هدف از این بخش، جمع‌بندی بی‌طرفانهٔ نتایج و بررسی نقاط قوت و ضعف هر رویکرد است.

\subsection{مقایسهٔ عملکرد در هدایت نقطه ثابت (ورودی پله)}
بررسی پاسخ سیستم به خطای موقعیت ناگهانی نشان داد که تنظیم‌گر خطی درجه دوم (\lr{LQR}) وابستگی شدیدی به اعتبار فرض‌های خطی‌سازی دارد. این کنترل‌کننده در مسافت‌های کوتاه ($0.5$ متر) عملکردی قابل‌قبول داشت، اما در جابجایی‌های بزرگ ($5$ متر) به دلیل درخواست نیروهای غیرواقع‌بینانه، منجر به خروج پاندول از محدودهٔ خطی و بروز ناپایداری فیزیکی شد.

در مقابل، کنترل‌کنندهٔ پیش‌بین مدل غیرخطی (\lr{NMPC}) با بهره‌گیری از معادلات دقیق دینامیکی و فعال‌سازی مکانیزم تغییر طول کابل، توانست بر این چالش غلبه کند و پایداری کالسکه را در جابجایی بلندبرد $5$ متری نیز حفظ نماید. با این وجود، تحلیل‌ها نشان داد که دستیابی به این پایداری سراسری، نیازمند یک مصالحهٔ مهندسی در تنظیم وزن‌های بهینه‌ساز است. برای جلوگیری از بروز ناپایداری در مسافت‌های طولانی، جریمهٔ سنگینی بر نرخ تغییرات عملگرها اعمال شد که این رویکردِ محافظه‌کارانه، در مسافت‌های کوتاه ($0.5$ متر) باعث کندی در ترمزگیری و بروز فراجهش (\lr{Overshoot}) در موقعیت کالسکه گردید.

\subsection{مقایسهٔ معماری‌های مبتنی بر برنامه‌ریز مسیر}
با ورود برنامه‌ریز مسیر سینماتیکی به حلقهٔ کنترل، ماهیت رفتار هر دو کنترل‌کننده دگرگون شد. تحلیل نتایج به‌دست‌آمده، یک واقعیت مهم مهندسی را آشکار می‌سازد: از منظر کیفیت ردیابی، همواری مسیر و عدم بروز فراجهش، معماریِ ترکیب‌یافته از \lr{LQR}، فیدفوروارد تحلیلی و برنامه‌ریز مسیر، عملکردی به مراتب روان‌تر و مطلوب‌تر از سیستم پیش‌بین غیرخطی ارائه داد. در این معماری، کالسکه بر روی پروفایل مرجع حرکت کرده و نوسانات آونگ نیز در بازهٔ بسیار کوچکی مهار شد.

با این حال، برتری عملکرد ردیابی در \lr{LQR} با یک محدودیت ساختاری همراه است. این کنترل‌کننده فاقد هوشمندی ریاضی برای درک محدودیت‌های سیستم است. به عبارت دیگر، با وجود اینکه \lr{LQR} در این شبیه‌سازی‌ها قید نوسان $15$ درجه را رعایت کرد، اما این موفقیت صرفاً ناشی از طراحی دقیق مسیر مرجع بود و هیچ‌گونه تضمین تحلیلی برای عدم نقض قیود (در صورت بروز اغتشاشات محیطی یا تغییرات جرمی بار) وجود ندارد. علاوه بر این، \lr{LQR} به دلیل ثابت فرض کردن طول کابل، از ظرفیت میرایی فعال سیستم بی‌بهره ماند.

در سوی دیگر، تلفیق \lr{NMPC} با برنامه‌ریز مسیر، اگرچه نیازمند توان پردازشی بسیار بالاتری است و ممکن است در برخی معیارهای ردیابیِ گذرا به نرمیِ رویکرد کلاسیکِ همراه با فیدفوروارد عمل نکند، اما دو مزیت منحصربه‌فرد را به سیستم می‌افزاید: نخست، توانایی استفاده از تغییرات طول کابل برای جذب انرژی نوسانی (\lr{Active Damping})، و دوم، تضمین صریح و ریاضیاتی برای ارضای قیود فیزیکی (مانند کران نوسان و محدودیت عملگرها) در هر لحظه از زمان.

\subsection{نتیجه‌گیری نهایی}
یافته‌های این پژوهش نشان می‌دهد که هیچ‌یک از معماری‌های بررسی‌شده، برتری مطلقی بر دیگری ندارند و انتخاب ساختار نهایی مستلزم شناخت دقیق نیازهای مسئله است. در کاربردهای صنعتی متعارف که مسیرها از پیش تعیین شده‌اند و اغتشاشات محیطی ناچیز است، استفاده از کنترل‌کنندهٔ کلاسیک به همراه فیدفوروارد و برنامه‌ریز مسیر، به دلیل سادگیِ پیاده‌سازی و کیفیت عالی در ردیابی، توجیه‌پذیرترین گزینه است. اما در محیط‌های پویاتر که رعایت دقیق مرزهای ایمنی (نظیر حفظ نوسان در یک بازهٔ اکید) حیاتی بوده و امکان استفاده از طول کابل برای میراسازی نوسان وجود دارد، استقرار کنترل‌کنندهٔ پیش‌بین غیرخطی به یک ضرورت تبدیل می‌شود.